\documentclass[11pt,a4paper]{article}

% ── Packages ──────────────────────────────────────────────────────────
\usepackage[utf8]{inputenc}
\usepackage[T1]{fontenc}
\usepackage{lmodern}
\usepackage[margin=2.5cm]{geometry}
\usepackage{booktabs}
\usepackage{tabularx}
\usepackage{array}
\usepackage{xcolor}
\usepackage{hyperref}
\usepackage{fancyhdr}
\usepackage{enumitem}
\usepackage{amssymb}
\usepackage{amsmath}
\usepackage{listings}
\usepackage{tcolorbox}

% ── Colours ───────────────────────────────────────────────────────────
\definecolor{ee8blue}{RGB}{0,90,160}
\definecolor{ee8gray}{RGB}{100,100,100}
\definecolor{notebg}{RGB}{235,245,255}
\definecolor{noteborder}{RGB}{70,130,200}
\definecolor{codebg}{RGB}{248,248,248}
\definecolor{keywd}{RGB}{0,90,160}
\definecolor{cmt}{RGB}{100,130,100}
\definecolor{strng}{RGB}{160,60,60}
\definecolor{machbg}{RGB}{245,240,255}
\definecolor{machborder}{RGB}{120,80,180}

% ── tcolorbox styles ─────────────────────────────────────────────────
\newtcolorbox{notebox}{
  colback=notebg, colframe=noteborder,
  fonttitle=\bfseries, title=Note,
  boxrule=0.5pt, aurc=2pt,
  left=6pt, right=6pt, top=4pt, bottom=4pt,
}

% ── Listings: C ──────────────────────────────────────────────────────
\lstdefinestyle{cstyle}{
  language=C,
  backgroundcolor=\color{codebg},
  basicstyle=\ttfamily\small,
  keywordstyle=\color{keywd}\bfseries,
  commentstyle=\color{cmt}\itshape,
  stringstyle=\color{strng},
  numbers=left,
  numberstyle=\tiny\color{ee8gray},
  numbersep=8pt,
  frame=single,
  rulecolor=\color{ee8gray!40},
  xleftmargin=1.5em,
  framexleftmargin=1.5em,
  tabsize=4,
  showstringspaces=false,
  breaklines=true,
}

% ── Listings: EE8 Assembly ───────────────────────────────────────────
\lstdefinelanguage{ee8asm}{
  morekeywords={LDI,ADD,SUB,AND,OR,XOR,SHL,SHR,MOV,IN,OUT,LT,EQ,JMP,JZ,JNZ,HALT},
  morecomment=[l]{;},
  sensitive=false,
}
\lstdefinestyle{asmstyle}{
  language=ee8asm,
  backgroundcolor=\color{codebg},
  basicstyle=\ttfamily\small,
  keywordstyle=\color{keywd}\bfseries,
  commentstyle=\color{cmt}\itshape,
  numbers=left,
  numberstyle=\tiny\color{ee8gray},
  numbersep=8pt,
  frame=single,
  rulecolor=\color{ee8gray!40},
  xleftmargin=1.5em,
  framexleftmargin=1.5em,
  tabsize=4,
  showstringspaces=false,
}

% ── Hyperref setup ────────────────────────────────────────────────────
\hypersetup{
  colorlinks=true,
  linkcolor=ee8blue,
  urlcolor=ee8blue,
  pdftitle={EE8 Count-to-5 Program},
  pdfauthor={EG2300},
}

% ── Header / footer ──────────────────────────────────────────────────
\pagestyle{fancy}
\fancyhf{}
\fancyhead[L]{\small\textcolor{ee8gray}{EG2300 --- EE8 CPU}}
\fancyhead[R]{\small\textcolor{ee8gray}{Count-to-5 Loop}}
\fancyfoot[C]{\small\textcolor{ee8gray}{\thepage}}
\renewcommand{\headrulewidth}{0.4pt}
\renewcommand{\footrulewidth}{0pt}

% ── Convenience ──────────────────────────────────────────────────────
\newcommand{\reg}[1]{\texttt{R#1}}
\newcommand{\hex}[1]{\texttt{0x#1}}

\begin{document}

% ── Title ────────────────────────────────────────────────────────────
\begin{center}
  {\LARGE\bfseries EE8 --- Count-to-5 Loop}\\[6pt]
  {\large C $\,\to\,$ Assembly $\,\to\,$ Machine Code}
\end{center}
\vspace{1em}

\noindent
A program that counts from 0 to 5 in \reg{2}, then resets and repeats
forever. Shown at three levels of abstraction. This is a basic test
program that exercises LDI, ADD, LT, JZ, and JMP --- no I/O required.

% ═════════════════════════════════════════════════════════════════════
\section{C (Equivalent Logic)}
% ═════════════════════════════════════════════════════════════════════

\begin{lstlisting}[style=cstyle]
#include <stdint.h>

int main(void)
{
    const uint8_t limit = 5;
    const uint8_t step  = 1;

    while (1) {
        uint8_t counter = 0;

        while (counter < limit)
            counter += step;
        /* counter == 5 here, then resets */
    }
}
\end{lstlisting}

\begin{notebox}
The EE8 has no stack, no function calls, and 8-bit unsigned data.
\texttt{uint8\_t} variables map directly to registers; the outer
\texttt{while(1)} mirrors the unconditional \texttt{JMP} back to the
comparison.
\end{notebox}

% ═════════════════════════════════════════════════════════════════════
\section{EE8 Assembly}
% ═════════════════════════════════════════════════════════════════════

\noindent
Register allocation: \reg{0}\,=\,limit (5),\quad \reg{1}\,=\,step (1),\quad
\reg{2}\,=\,counter,\quad \reg{3}\,=\,comparison result.

\begin{lstlisting}[style=asmstyle]
; ── Count-to-5 loop ──────────────────────────────────
; R2 counts 0, 1, 2, 3, 4, 5 then resets to 0.
;
0:      LDI  R0, 5          ; R0 = limit
1:      LDI  R1, 1          ; R1 = increment
2:      LDI  R2, 0          ; R2 = counter = 0
3:      LT   R2, R0         ; R3 = (R2 < R0) ? 1 : 0
4:      JZ   R3, 2          ; if R2 >= limit, restart
5:      ADD  R1, R2, R2     ; R2 = R2 + 1
6:      JMP  3              ; next iteration
\end{lstlisting}

% ═════════════════════════════════════════════════════════════════════
\section{EE8 Machine Code}
% ═════════════════════════════════════════════════════════════════════

\begin{center}
\begin{tabular}{>{\ttfamily}c >{\ttfamily}c l l}
  \toprule
  \textrm{\textbf{Addr}} & \textrm{\textbf{Hex}} &
  \textbf{Binary}        & \textbf{Instruction} \\
  \midrule
  0 & A005 & \texttt{1010\,0000\,0000\,0101} & \texttt{LDI R0, 5}  \\
  1 & A401 & \texttt{1010\,0100\,0000\,0001} & \texttt{LDI R1, 1}  \\
  2 & A800 & \texttt{1010\,1000\,0000\,0000} & \texttt{LDI R2, 0}  \\
  3 & 8800 & \texttt{1000\,1000\,0000\,0000} & \texttt{LT\;\;R2, R0} \\
  4 & CC02 & \texttt{1100\,1100\,0000\,0010} & \texttt{JZ\;\;R3, 2}  \\
  5 & 0680 & \texttt{0000\,0110\,1000\,0000} & \texttt{ADD R1, R2, R2} \\
  6 & B003 & \texttt{1011\,0000\,0000\,0011} & \texttt{JMP 3}       \\
  \midrule
  7--F & 0000 & \texttt{0000\,0000\,0000\,0000} & \textrm{(unused)} \\
  \bottomrule
\end{tabular}
\end{center}

\begin{notebox}
\textbf{Encoding notes:}
I-type (LDI) places Rd at bits [11:10]. R-type uses Rs1\,[11:10],
Rs2\,[9:8], Rd\,[7:6]. J-type uses Rcond\,[11:10], address\,[3:0].

\medskip
\textbf{LT encoding:}
\texttt{LT R2, R0} has Rs1=R2 at [11:10]=10, Rs2=R0 at [9:8]=00.
The result always goes to R3 (hardwired by \texttt{flag\_write}).
\end{notebox}

% ═════════════════════════════════════════════════════════════════════
\section{Execution Trace}
% ═════════════════════════════════════════════════════════════════════

\begin{center}
\begin{tabular}{c >{\ttfamily}c >{\ttfamily}c >{\ttfamily}c >{\ttfamily}c >{\ttfamily}c l}
  \toprule
  \textbf{Cycle} & \textrm{\textbf{PC}} & \textrm{\textbf{R0}} &
  \textrm{\textbf{R1}} & \textrm{\textbf{R2}} & \textrm{\textbf{R3}} & \textbf{Action} \\
  \midrule
  1  & 0 & 05 & {-}{-} & {-}{-} & {-}{-} & \texttt{LDI R0, 5} \\
  2  & 1 & 05 & 01 & {-}{-} & {-}{-} & \texttt{LDI R1, 1} \\
  3  & 2 & 05 & 01 & 00 & {-}{-} & \texttt{LDI R2, 0} \\
  \midrule
  4  & 3 & 05 & 01 & 00 & 01 & \texttt{LT}: $0 < 5 \to$ R3=1 \\
  5  & 4 & 05 & 01 & 00 & 01 & \texttt{JZ} not taken (R3$\neq$0) \\
  6  & 5 & 05 & 01 & 01 & 01 & \texttt{ADD}: R2$\leftarrow$1 \\
  7  & 6 & 05 & 01 & 01 & 01 & \texttt{JMP 3} \\
  \midrule
  8  & 3 & 05 & 01 & 01 & 01 & \texttt{LT}: $1 < 5 \to$ R3=1 \\
  9  & 4 & 05 & 01 & 01 & 01 & \texttt{JZ} not taken \\
  10 & 5 & 05 & 01 & 02 & 01 & R2$\leftarrow$2 \\
  11 & 6 & 05 & 01 & 02 & 01 & \texttt{JMP 3} \\
  \midrule
  \multicolumn{7}{c}{$\vdots$ \textrm{(R2 increments: 3, 4)}} \\
  \midrule
  20 & 3 & 05 & 01 & 05 & 00 & \texttt{LT}: $5 < 5 \to$ R3=0 \\
  21 & 4 & 05 & 01 & 05 & 00 & \texttt{JZ} taken $\to$ addr 2 \\
  22 & 2 & 05 & 01 & 00 & 00 & \texttt{LDI R2, 0} (reset) \\
  \bottomrule
\end{tabular}
\end{center}

\noindent
The loop runs 5 iterations per cycle. R2 counts 0\,$\to$\,1\,$\to$\,2\,$\to$\,3\,$\to$\,4\,$\to$\,5, then the comparison $5 < 5$ fails (R3=0), \texttt{JZ} branches back to address~2, and R2 resets to~0.

\end{document}
